\subsection{Scaling across datasets}\label{sec:disc:scaling}

The human dataset, six times larger than primates14.chr6, scales much
better, which fits the explanation in Section~\ref{sec:disc:small}: a
contig pair is never split, so the most expensive pair bounds the
alignment stage. hsapiens90.chr6 has 1,422 contigs against 82, and its most
expensive pair is about 0.5\% of the alignment work rather than 7\%. Its
work could in principle be spread over roughly 180 ranks before that pair
sets the stage time, against about 14 for primates14.chr6. How many contigs
an input has, and how evenly its work falls on them, may therefore matter
more for scaling than its size.

The task bound does not explain the remaining imbalance on the human data.
At 64 ranks its slowest rank aligns for 149\,s, more than four times its
most expensive pair (34\,s) and 1.49 times the average (100\,s). Since no
single task is that large, the gap most likely comes from the static
assignment by length product, which does not track alignment cost (Section~\ref{sec:disc:future}). A perfectly balanced
assignment would shorten the alignment stage by at most this factor of
about 1.5 at 64 ranks.
